% =====================================================================
%  CollusionGraph — Internal Team Report
%  Plain-language, complete walkthrough of the whole project.
%
%  Build:  pdflatex collusiongraph_internal_report.tex   (run twice for
%          the table of contents and cross-references), or
%          latexmk -pdf collusiongraph_internal_report.tex
%
%  Only standard TeX Live / MiKTeX packages are used — no downloads.
%
%  HOUSE RULE: every number in this report is copied from a measured
%  artifact recorded in PROGRESS.md, docs/model_card.md or
%  docs/reproducibility.md. Nothing here is estimated or remembered.
% =====================================================================
\documentclass[11pt,a4paper]{article}

\usepackage[T1]{fontenc}
\usepackage[utf8]{inputenc}
\usepackage[margin=2.4cm]{geometry}
\usepackage{lmodern}
\usepackage{microtype}
\usepackage{booktabs}
\usepackage{longtable}
\usepackage{array}
\usepackage{xcolor}
\usepackage{graphicx}
\usepackage{amsmath}
\usepackage{amssymb}
\usepackage{enumitem}
\usepackage{tcolorbox}
\usepackage{fancyhdr}
\usepackage{titlesec}
\usepackage[hidelinks,bookmarks=true]{hyperref}

\definecolor{ink}{HTML}{1A1C20}
\definecolor{accent}{HTML}{0E7490}
\definecolor{soft}{HTML}{F4F5F7}
\definecolor{rule}{HTML}{D9DCE1}
\definecolor{warn}{HTML}{B45309}

\hypersetup{colorlinks=true,linkcolor=accent,urlcolor=accent}

\titleformat{\section}{\Large\bfseries\color{ink}}{\thesection}{0.7em}{}
\titleformat{\subsection}{\large\bfseries\color{ink}}{\thesubsection}{0.6em}{}

\pagestyle{fancy}
\fancyhf{}
\lhead{\small\textcolor{ink}{CollusionGraph — Internal Team Report}}
\rhead{\small\thepage}
\renewcommand{\headrulewidth}{0.4pt}

% --- helper boxes -----------------------------------------------------
\newtcolorbox{plain}[1][]{
  colback=soft, colframe=rule, boxrule=0.5pt, arc=2pt,
  left=8pt, right=8pt, top=6pt, bottom=6pt, #1}

\newtcolorbox{keybox}[1][]{
  colback=white, colframe=accent, boxrule=1pt, arc=2pt,
  left=8pt, right=8pt, top=6pt, bottom=6pt, #1}

\newtcolorbox{warnbox}[1][]{
  colback=white, colframe=warn, boxrule=1pt, arc=2pt,
  left=8pt, right=8pt, top=6pt, bottom=6pt, #1}

% "in plain words" lead-in
\newcommand{\pw}[1]{\textit{\textcolor{accent}{In plain words:}} #1}
\newcommand{\code}[1]{\texttt{\small #1}}

\setlist[itemize]{topsep=3pt,itemsep=2pt,parsep=0pt,leftmargin=1.4em}
\setlist[enumerate]{topsep=3pt,itemsep=2pt,parsep=0pt,leftmargin=1.6em}

% =====================================================================
\begin{document}

\begin{titlepage}
\centering
\vspace*{2.5cm}
{\Huge\bfseries CollusionGraph\par}
\vspace{0.6cm}
{\LARGE Finding Hidden Groups That Cheat Together\par}
\vspace{1.4cm}
{\Large\bfseries Internal Team Report\par}
\vspace{0.4cm}
{\large A complete, plain-language walkthrough of the whole project\par}
\vspace{2.2cm}

\begin{plain}
\raggedright
\textbf{Who this is for.} Everyone on the team — including people who have
never trained a model. It also answers the questions an experienced AI
researcher or examiner would ask.

\vspace{4pt}
\textbf{How it is written.} Ordinary English first, technical name second.
Every specialised word is defined in the Glossary (Section~\ref{sec:glossary}).
If you read only one section, read Section~\ref{sec:story}.

\vspace{4pt}
\textbf{Honesty rule.} Every number here was measured by our own test harness
and is recorded in the project ledger (\code{PROGRESS.md}), the model card
(\code{docs/model\_card.md}) or the reproducibility map
(\code{docs/reproducibility.md}). Results that make the project look
\emph{worse} are included, because a report that hides them is useless.
\end{plain}

\vfill
{\small Capstone project · Research prototype (TRL 3--4) · Not a production system\par}
{\small Screening signal only --- never a determination of guilt\par}
\end{titlepage}

\tableofcontents
\newpage

% =====================================================================
\section{The Project in One Page}
\label{sec:story}

\subsection{The problem}

Two very different crimes have the same hidden shape.

\begin{itemize}
  \item \textbf{Money laundering.} Criminals move dirty money through many
        bank accounts or crypto wallets so it looks clean. No single payment
        looks wrong. The \emph{pattern of payments} is what is wrong.
  \item \textbf{Bid rigging (cartels).} Companies that are supposed to compete
        for government contracts secretly agree who wins. Each bid looks
        normal. The \emph{pattern of who bids against whom, and who wins in
        turn}, is what is wrong.
\end{itemize}

\begin{keybox}
\textbf{The core idea of this project.} In both crimes the guilty parties are
not lone individuals --- they are a \emph{group} connected to each other.
If you draw the data as a network of dots and lines, the crime shows up as a
\emph{shape} in the network. So instead of asking ``is this one payment
suspicious?'' we ask ``is this whole group of connected entities behaving like
a coordinated ring?''
\end{keybox}

\subsection{What we built}

One machine-learning system that reads both kinds of data, turns them into the
same internal format, and produces a short, ranked list of suspicious
\emph{groups} --- each with an explanation of \emph{why} it was flagged. On top
of that sits a web dashboard where an investigator can look through the list,
inspect the network, read the evidence, and ask an AI assistant questions.

\subsection{The five things that make it different}

\begin{enumerate}
  \item \textbf{One system, two crimes.} About 70\% of the code is shared;
        only the data readers and the red-flag vocabularies differ.
  \item \textbf{It ranks groups, not single items.} The output is a queue of
        communities, because that is what an investigator actually opens.
  \item \textbf{Every alert carries an explanation} --- the exact sub-network
        that triggered it, the pattern name, and a citation to an official
        red-flag list (FATF for money laundering, OECD for bid rigging).
  \item \textbf{Fixed review budget.} We never report ``accuracy''. We report
        ``if your team can review 100 cases this week, how many of those 100
        are real?'' That is the number that matters operationally.
  \item \textbf{Honest reporting.} Where a simple method beats our deep-learning
        model, we say so and publish the gap (see Section~\ref{sec:results}).
\end{enumerate}

% =====================================================================
\section{Glossary --- Every Term We Use}
\label{sec:glossary}

Read this once and the rest of the report will be easy. Terms are grouped, and
within each group the simplest come first.

\subsection{Network / graph words}

\begin{longtable}{@{}p{4.1cm}p{11.4cm}@{}}
\toprule
\textbf{Term} & \textbf{Meaning in plain words} \\
\midrule
\endhead
\textbf{Graph} (network) & A picture made of dots joined by lines. Dots =
things, lines = relationships. Everything in this project is a graph. \\
\textbf{Node} (vertex) & One dot. A bank account, a transaction, a company, a
tender. \\
\textbf{Edge} & One line joining two dots. A payment, an award, a bid. \\
\textbf{Directed edge} & A line with an arrow. Money goes \emph{from} A
\emph{to} B --- the direction matters, because money flowing in is a different
crime from money flowing out. \\
\textbf{Degree} & How many lines touch a dot. ``In-degree'' = arrows coming in,
``out-degree'' = arrows going out. A shop has high in-degree; a payroll account
has high out-degree. \\
\textbf{Neighbour} & A dot directly joined to another dot by one line. \\
\textbf{Subgraph} & A small piece cut out of a big graph --- e.g.\ one
suspicious group and the lines between them. \\
\textbf{Ego-network} & The subgraph made of one chosen dot plus everything
within a few steps of it. What you see in the Graph Explorer screen. \\
\textbf{Motif} & A small, repeating \emph{shape} in the network that has a
meaning. Our nine motifs are listed in Table~\ref{tab:motifs}. \\
\textbf{Community} & A cluster of dots that are much more connected to each
other than to the rest of the network. Our alerts are communities. \\
\textbf{Clique} & A group where \emph{everyone} is connected to
\emph{everyone else} --- the densest possible community. \\
\textbf{k-core} & A measure of how deep inside a dense region a dot sits. High
k-core means ``buried in a tight cluster''. \\
\textbf{Clustering coefficient} & How often your neighbours are also each
other's neighbours. High = your contacts all know each other (cliquey). \\
\textbf{Bipartite graph} & A graph with two kinds of dots where lines only go
between kinds, never within. Firms and tenders: a firm bids on a tender, but a
firm never ``bids on'' another firm. \\
\textbf{Projection} & Collapsing a bipartite graph into one kind of dot. From
``firms bid on tenders'' we make ``firm A and firm B bid on the same tender'' ---
a firm-to-firm graph. This is how we get \textbf{co-bid} links. \\
\bottomrule
\end{longtable}

\subsection{Machine-learning words}

\begin{longtable}{@{}p{4.1cm}p{11.4cm}@{}}
\toprule
\textbf{Term} & \textbf{Meaning in plain words} \\
\midrule
\endhead
\textbf{Model} & A program that learns patterns from examples instead of being
given fixed rules. \\
\textbf{Feature} & One measurable fact about a node, written as a number.
``In-degree = 14'' is a feature. Models only eat numbers. \\
\textbf{Label} & The right answer for an example: illicit, licit, or unknown. \\
\textbf{Training} & Showing the model labelled examples so it adjusts itself to
predict better. \\
\textbf{Validation} & A held-back slice used \emph{during} development to decide
when to stop training and which settings to keep. \\
\textbf{Test} & A slice the model never sees until the very end. The only
honest measure of quality. \\
\textbf{Supervised} & Learning from labelled examples (``these were cartels''). \\
\textbf{Unsupervised} & Learning with \emph{no} labels --- the model just learns
what ``normal'' looks like and flags whatever does not fit. \\
\textbf{Epoch} & One full pass through the training data. \\
\textbf{Early stopping} & Stop training when the validation score stops
improving, so the model does not over-memorise. \\
\textbf{Overfitting} & When a model memorises the training data instead of
learning the general pattern --- great on training, useless on new data. \\
\textbf{Seed} & The starting random number. Different seeds give slightly
different models. We always run 5 seeds and report the average \emph{and} the
spread, because a single seed can flatter you. \\
\textbf{Hyperparameter} & A setting you choose rather than learn: how many
layers, the learning rate, the loss function. \\
\textbf{Embedding} & A list of numbers that summarises a node's position and
behaviour in the network, produced by the model. \\
\textbf{Inference / scoring} & Using a trained model to produce scores for new
data. \\
\textbf{Ablation} & Deliberately removing one part of the system to measure how
much that part actually contributed. \\
\textbf{Baseline} & A deliberately simple method used as a yardstick. If the
complicated model cannot beat it, the complication is not earning its place. \\
\bottomrule
\end{longtable}

\subsection{Words specific to this project}

\begin{longtable}{@{}p{4.1cm}p{11.4cm}@{}}
\toprule
\textbf{Term} & \textbf{Meaning in plain words} \\
\midrule
\endhead
\textbf{IR} (Intermediate Representation) & Our one common data format. Every
dataset, financial or procurement, is converted into the same four files, so
everything downstream works on both domains without changes. \\
\textbf{Adapter} & The small piece of code that converts one raw dataset into
the IR. \\
\textbf{Alert} & One suspicious \emph{community}, with a rank, a risk score, and
an explanation bundle. The unit an investigator reviews. \\
\textbf{Alert queue} & The ranked list of alerts, cut at a review budget. \\
\textbf{Budget (k)} & How many cases the team can actually review. All our
headline numbers are stated at a budget. \\
\textbf{Explanation bundle} & A structured JSON file attached to every alert:
the minimal sub-network, the motif name, the evidence, the red-flag citation,
fidelity scores, and the locked ethics caveat. \\
\textbf{Motif matcher} & Transparent, rule-based code that names the shape it
finds (cycle, rotation, cover bid \ldots). It does not use the neural network,
so it is an independent check on it. \\
\textbf{Red flag} & An officially published warning sign. We cite FATF
indicators (money laundering) and OECD indicators (bid rigging). \\
\textbf{Floor} & The simplest possible scorer (structural z-scores). It exists
so we always know what ``no cleverness at all'' achieves. \\
\textbf{Calibration} & Rescaling each model's raw output so the numbers mean the
same thing before we combine them. Without this, combining models destroys
their quality (we measured it --- Section~\ref{sec:results}). \\
\textbf{Ensemble / fusion} & Combining several models' scores into one. \\
\textbf{Roll-up} & Turning per-node scores into a per-community score. \\
\textbf{NMS} (non-maximum suppression) & Removing near-duplicate alerts, so one
big group does not fill the whole queue in ten slightly different versions. \\
\textbf{Leakage} & When information from the test period accidentally reaches
the model during training. It inflates scores and invalidates results. We have
a dedicated test suite that runs in CI to prevent it. \\
\textbf{Strict-inductive split} & Our anti-leakage rule: at training time the
model may not even \emph{see the connections} that belong to the test period,
not just the labels. \\
\textbf{As-of discipline} & Every feature is computed using only information
that existed at that moment in time --- never the future. \\
\textbf{Injection study} & We plant synthetic, known-fake cartels into a real
network and measure how many we can recover. Because we planted them, we know
the perfect answer. \\
\textbf{Copilot} & The AI chat assistant inside the dashboard. \\
\textbf{Grounding} & The rule that the Copilot may only state numbers that
actually appear in the data it retrieved. \\
\textbf{Artifact} & An output file produced by a run (scores, alerts,
explanations, metrics). The website only ever reads artifacts. \\
\textbf{Trust boundary} & The line between the part that can change results
(training) and the part that only reads them (the website). \\
\bottomrule
\end{longtable}

\subsection{Measurement words}

\begin{longtable}{@{}p{4.1cm}p{11.4cm}@{}}
\toprule
\textbf{Term} & \textbf{Meaning in plain words} \\
\midrule
\endhead
\textbf{Prevalence} & How common the bad thing is. If 2\% of transactions are
illicit, prevalence is 0.02. This is the ``do nothing clever'' score. \\
\textbf{Precision@k} & Out of the top $k$ alerts, what fraction are real? If
Precision@100 = 0.81, then 81 of your 100 reviewed cases were genuine. \\
\textbf{Recall@k} & Out of all the bad things that exist, what fraction did the
top $k$ catch? \\
\textbf{AUC-PR} & A single summary score for ranking quality across all
budgets. \textbf{Crucially, its ``random guessing'' value equals the
prevalence}, so AUC-PR $=0.47$ at prevalence $0.065$ means roughly
7$\times$ better than chance. We always quote prevalence beside it. \\
\textbf{Lift} & Score divided by prevalence. ``How many times better than
guessing.'' \\
\textbf{Why not accuracy?} & At 2\% prevalence, a model that says ``nothing is
ever bad'' is 98\% accurate and completely useless. We never headline accuracy. \\
\textbf{Mean $\pm$ std} & Average across 5 seeds, plus the spread. A big spread
means the result is unstable and must not be quoted as a single number. \\
\textbf{Confidence interval (CI)} & A range we are 95\% sure the true value sits
in. \\
\textbf{Paired bootstrap} & A statistical test: re-sample the same test cases
thousands of times and see whether model A really beats model B, or whether the
gap is just luck. \\
\textbf{Fidelity} & Does the explanation actually matter? Remove the highlighted
sub-network --- if the model's score collapses, the explanation was real. \\
\bottomrule
\end{longtable}

% =====================================================================
\section{The Nine Patterns We Look For}

Every alert is named with one of these shapes. Five come from the money-crime
side, five from the procurement side, and \emph{hidden common control} appears
in both --- which is one of the project's nicest findings: the same shape means
the same thing in both worlds.

\begin{table}[h]
\centering
\small
\begin{tabular}{@{}p{2.6cm}p{5.9cm}p{5.9cm}@{}}
\toprule
\textbf{Shape} & \textbf{In money laundering} & \textbf{In bid rigging} \\
\midrule
Circular & \textbf{cycle} --- money leaves an account and comes back through
intermediaries (round-tripping). &
\textbf{rotation} --- the same firms take turns winning. \\
\addlinespace
Converging & \textbf{fan\_in} --- many small deposits, all just under the
reporting threshold, land in one account (``smurfing''). &
\textbf{cover\_bid} --- losing bids sit implausibly just above the winner, to
make a fixed price look competitive. \\
\addlinespace
Spreading out & \textbf{fan\_out} --- money is sprayed to many accounts at once
to break the trail (``layering''). &
\textbf{partition} --- firms carve up the market: you take this region, I take
that one, and we never bid against each other. \\
\addlinespace
Hidden control & \textbf{common\_control} --- accounts that look unrelated share
an owner, agent or address. &
\textbf{common\_control} --- ``rival'' firms share directors or addresses. \\
\addlinespace
Tight cluster & \textbf{pass\_through} --- chains where money arrives and leaves
almost immediately, keeping nothing. &
\textbf{clique} (coordinated cluster) --- the same firms always bid together,
with suspiciously similar prices. \\
\bottomrule
\end{tabular}
\caption{The nine motif families. \code{common\_control} is shared, so the
system genuinely spans both domains.}
\label{tab:motifs}
\end{table}

\begin{plain}
\pw{A cycle is money doing a lap and coming home. A fan-in is a funnel. A
fan-out is a sprinkler. Common control is two ``strangers'' who share a phone
number. A clique is a group where everybody knows everybody.}
\end{plain}

% =====================================================================
\section{The Data}
\label{sec:data}

\subsection{Which datasets, and why each one}

We use six public datasets. Nothing private, nothing personal, nothing from any
UAE institution. Raw data is \emph{never} committed to the repository; we ship
download scripts plus checksums so anyone can fetch and verify the exact same
files.

\begin{longtable}{@{}p{1.1cm}p{3.5cm}p{4.2cm}p{5.8cm}@{}}
\toprule
\textbf{ID} & \textbf{Dataset} & \textbf{What it is} & \textbf{Why we need it} \\
\midrule
\endhead
D1 & Elliptic / Elliptic++ & Real Bitcoin transactions, labelled illicit/licit &
Our main \emph{financial} dataset. Real crime, real labels, hard imbalance. \\
D2 & IBM AMLworld HI-Small & Synthetic bank transactions with perfect ground
truth for 8 laundering patterns & Lets us check the detector on patterns whose
answer we know exactly; gives real money amounts. \\
D3 & García Rodríguez et al. & Auctions from Brazil, Italy, Japan,
Switzerland, USA --- collusive vs.\ competitive & Has \emph{losing} bids, so
real co-bid networks and price screens work. Used for cross-market transfer. \\
D4 & Mendeley EU cartel & EU public contracts with 73 confirmed cartel cases &
Our main \emph{procurement} dataset with real labels. \\
D5 & OCDS --- Georgia OpenTender & 451{,}346 real public-tender records,
\emph{no labels} & The large realistic playground where we plant synthetic
cartels and measure recovery at scale. \\
\bottomrule
\end{longtable}

\subsection{The spread of the data --- exact measured numbers}

These were verified by our own exploratory analysis, not copied from papers.

\begin{longtable}{@{}p{4.3cm}p{11.2cm}@{}}
\toprule
\textbf{Dataset} & \textbf{Measured shape} \\
\midrule
\endhead
\textbf{Elliptic (base)} &
203{,}769 transaction nodes; 234{,}355 directed edges; 49 time steps;
166 features per node.
Labels: 4{,}545 illicit ($\approx$2\%), 42{,}019 licit ($\approx$21\%),
\textbf{the remaining $\approx$77\% unknown}. \\
\addlinespace
\textbf{Elliptic++ (extension)} &
Adds $\approx$822{,}942 wallet/actor nodes with 56 named features and 1.27M+
temporal interactions. 183 features on the transaction view. \\
\addlinespace
\textbf{AMLworld HI-Small} &
515{,}088 accounts; 5{,}078{,}345 payment edges with real amounts; complete
ground truth for 8 pattern types. \\
\addlinespace
\textbf{Mendeley EU} &
15{,}616 contract rows; \textbf{73 distinct cartel cases} (matches the source
exactly); \code{is\_cartel=1} on 6{,}548 rows (41.9\%);
7 anonymised countries; 2004--2021.
Per-country rows range 162 to 7{,}860 and per-country cartel share ranges
\textbf{9.9\% to 59.6\%} --- the countries are very unequal, which matters for
transfer tests.
Graph: 14{,}555 nodes / 24{,}251 edges. \\
\addlinespace
\textbf{García Rodríguez} &
77{,}007 nodes / 111{,}046 edges; firm identities available on 4 of 6 markets
(Japan, Italy, Brazil, USA carry competitor lists; the two Swiss markets do
not). \\
\addlinespace
\textbf{OCDS Georgia} &
451{,}346 compiled releases (2010--2025) $\rightarrow$
\textbf{488{,}300 nodes / 1{,}449{,}077 edges}, including
\textbf{687{,}336 \code{bids\_on} edges with identified losing bidders}.
Roughly 50$\times$ the node count of Mendeley. Zero undated releases, zero
id-less bids. \\
\bottomrule
\end{longtable}

\begin{warnbox}
\textbf{A caveat we must never hide (and an examiner will ask about).}
The Mendeley dataset is a \textbf{case-control sample}, not a population sample.
Its 41.9\% cartel rate is an artifact of how the researchers assembled it --- the
real world is nowhere near 42\% cartels. Therefore \emph{absolute} precision
numbers on Mendeley are not real-world rates; only \emph{lift-style} comparisons
(``how much better than the 0.358 prevalence baseline'') are meaningful, and only
comparisons within the same fold are fair. This is stated in the dataset
datasheet, the model card and the paper.
\end{warnbox}

\subsection{The class-imbalance problem}

In Elliptic, about 2\% of transactions are illicit. This is the central
technical difficulty:

\begin{itemize}
  \item A model can score 98\% ``accuracy'' by never flagging anything.
  \item The model sees 50 normal examples for every bad one, so it is tempted to
        ignore the bad class entirely.
  \item Ordinary metrics (accuracy, ROC-AUC) look great while the system is
        useless.
\end{itemize}

Our four answers: (1) \textbf{focal loss}, which makes the model pay more
attention to the rare, hard examples; (2) \textbf{budget-first metrics}
(Precision@k) instead of accuracy; (3) \textbf{AUC-PR always quoted against
prevalence}; (4) \textbf{synthetic injection} to create extra known-bad
sub-networks without distorting the real ones.

\subsection{How the data is split for training and testing}

This is the part that protects the credibility of every number in this report.

\begin{keybox}
\textbf{Rule 1 --- split by time, never randomly.} Crime data is a timeline.
If you shuffle it randomly, the model learns from the future to predict the
past, which is impossible in real life and inflates every score.
\end{keybox}

\begin{longtable}{@{}p{3.4cm}p{12.1cm}@{}}
\toprule
\textbf{Dataset} & \textbf{Split actually used} \\
\midrule
\endhead
\textbf{Elliptic++} & Train on time steps 1--34, test on 35--49 (the standard
convention in the literature, so our numbers are comparable). Validation is a
\emph{temporal tail} of the training period --- never a random slice. \\
\textbf{Mendeley} & Within-country: train on contracts up to 2013, test from
2014 onwards. Across countries: \textbf{LOCO}. \\
\textbf{García} & \textbf{LOMO} --- leave one market out. \\
\textbf{OCDS Georgia} & Test window 2020--2024; no labels exist, so this is
used for injection studies only. \\
\bottomrule
\end{longtable}

\textbf{LOCO / LOMO} mean \emph{leave-one-country-out} / \emph{leave-one-market-out}:
train on six countries, test on the seventh, then rotate. It answers the
question a buyer actually cares about --- ``will this work in a country where you
have never seen a single labelled cartel?''

\begin{keybox}
\textbf{Rule 2 --- strict-inductive splitting.} It is not enough to hide the
test \emph{labels}. We hide the test-period \emph{connections} too. During
training the model literally cannot see edges that belong to the test period.

\pw{If you are predicting who will be friends next year, it is cheating to
already look at next year's friendship lines --- even if you cover up the names.}
Recent published work shows a gap of about 39.5 F1 points comes purely from this
kind of exposure. We verified our splitter withholds exactly 77{,}512
test-period edges on Elliptic++.
\end{keybox}

\begin{keybox}
\textbf{Rule 3 --- as-of features.} Every feature is computed using only data
that existed at that time. Verified on real data: the ``as-of step 34'' view
shows 136{,}265 nodes and is \emph{exactly equal} to the truncated graph.
\end{keybox}

\begin{keybox}
\textbf{Rule 4 --- leakage tests run automatically.} A 25-test suite with
negative controls runs on every push. Any change that reintroduces leakage
fails the build.
\end{keybox}

\subsection{Data cleaning and preparation}

We deliberately did \textbf{not} ``clean'' the data in the usual sense --- no
dropping of awkward rows, no filling in of missing values with averages. In
fraud detection the weird rows are often the interesting ones, and inventing
values would manufacture evidence. What we actually did:

\begin{enumerate}
  \item \textbf{Verify before use.} Every file is checksummed (SHA-256) against
        a committed manifest. The Mendeley hash matches the publisher's own
        official API hash.
  \item \textbf{Convert, don't alter.} Each adapter maps raw columns into the
        IR. Values are carried across unchanged; anything not in the schema is
        preserved in a \code{raw\_attrs} JSON field rather than discarded.
  \item \textbf{Handle missing structure honestly.} Mendeley rows with no buyer
        simply produce no \code{buys\_from} edge --- we do not invent a buyer.
        Where losing-bidder data does not exist, the co-bid features are
        \emph{switched off} for that country-year rather than filled with zeros
        (a zero would be a lie: it means ``no co-bidding'', not ``unknown'').
  \item \textbf{Keep unknown labels in the graph.} 77\% of Elliptic nodes are
        unlabelled. We do \emph{not} delete them. They do not contribute to the
        loss, but they \emph{do} carry messages, because they are the structure
        that defines the crime. This is an explicit, documented policy.
  \item \textbf{Missing numbers stay missing.} XGBoost handles NaN natively;
        for isolated nodes the neighbour-average features are NaN, not 0.
  \item \textbf{Normalise per graph, frozen at train time.} Features are
        z-scored (rescaled to mean 0, spread 1). \textbf{The scaling constants
        are computed on training data only and then frozen.} This mattered
        enormously --- see the audit in Section~\ref{sec:tuning}.
  \item \textbf{Deduplicate alerts, not data.} Overlap removal happens at the
        alert stage (NMS), where it is a reporting decision, not a data
        decision.
\end{enumerate}

% =====================================================================
\section{The Common Data Format (IR)}

Every dataset becomes four tables. After this point, no code cares which domain
it is looking at --- this is what makes ``one stack, two crimes'' real rather
than a slogan.

\begin{plain}
\code{nodes.parquet} \quad node\_id · node\_type · domain · time\_first\_seen ·
features · raw\_attrs\\
\code{edges.parquet} \quad src · dst · edge\_type · timestamp · amount ·
directed · raw\_attrs\\
\code{labels.parquet} \quad node\_id · label $\in$ \{illicit, licit, unknown\} ·
label\_source · confidence\\
\code{communities.parquet} \quad community\_id · member node\_ids · method
\end{plain}

\begin{longtable}{@{}p{3.2cm}p{5.9cm}p{5.9cm}@{}}
\toprule
 & \textbf{Financial} & \textbf{Procurement} \\
\midrule
\endhead
Node types & account / address, transaction & firm, tender, bid, lot, buyer \\
Edge types & \code{pays} (directed, timed, amount), \code{linked\_to} &
\code{awarded}, \code{buys\_from} (always available); \code{bids\_on},
\code{co\_bid}, \code{linked\_to} (where data allows) \\
\bottomrule
\end{longtable}

\textbf{Two design rules that came out of an internal audit:}

\begin{enumerate}
  \item \textbf{Award-network-first.} Because losing bids are missing in many
        countries, every procurement component must work using \emph{awards
        alone}; richer co-bid features activate only where the data exists. This
        prevents building something that quietly only works on one country.
  \item \textbf{Shared structural template.} Degrees, motif counts, clustering,
        k-core, burstiness, community statistics --- computed identically for
        both domains and z-scored per graph. \textbf{This shared channel is what
        makes cross-domain transfer technically possible at all}, since the raw
        features of Bitcoin and EU tenders have nothing in common.
\end{enumerate}

% =====================================================================
\section{The Algorithms --- What They Are and Why We Chose Them}
\label{sec:algos}

\subsection{Graph Neural Networks (GNNs) --- the heart}

\pw{A GNN is a system where every dot repeatedly asks its neighbours ``what are
you?'', summarises the answers, and updates its own description. After two
rounds, each dot's description contains information about its neighbours and its
neighbours' neighbours. So the dot is described by \emph{the company it keeps},
which is exactly what we need --- because a laundering account looks normal alone
and suspicious in context.}

Formally, each node has a vector; at each layer it aggregates its neighbours'
vectors, combines that with its own, and passes it through a small neural
network. We use two layers (two hops of context).

\subsubsection*{GATv2 --- attention (our main financial model)}

\pw{Not all neighbours matter equally. Attention lets each node \emph{learn how
much to listen to} each neighbour. A node with 200 neighbours where 3 matter can
now focus on those 3.}
\textbf{Bonus:} the attention weights are readable, so they double as an
explanation of which links drove a score.

\textbf{Why GATv2 and not the original GAT?} The original GAT has a proven
mathematical flaw called \emph{static attention}: the ranking of which
neighbours are important ends up the same for every node, regardless of that
node. GATv2 fixes this at identical computational cost. Choosing the older one
would have been a strictly worse decision for no saving.

\subsubsection*{GraphSAGE --- the scalable workhorse}

\pw{Instead of looking at all neighbours (impossible when a node has 100{,}000),
sample a fixed number. This makes it work on huge graphs and, importantly, on
\emph{nodes it has never seen before}.}
We use a direction-aware version with separate aggregations for incoming and
outgoing edges.

\subsubsection*{R-GCN --- relation types (our main procurement model)}

\pw{In procurement, ``firm was awarded tender'' and ``firm bid on tender'' and
``firm shares a director with firm'' are completely different relationships.
R-GCN learns a separate rule for each type of line instead of treating all lines
as the same.}

\subsubsection*{Direction handling --- a small decision with a big effect}

Money flow and awards are directed, and \emph{direction is the crime}: fan-in
(funnel) and fan-out (sprinkler) are opposite offences. So every model receives
\textbf{both directions plus a flag saying which is which}. Our ablation shows
this is the single most important component in the whole stack
(Section~\ref{sec:results}).

\subsection{Autoencoders --- finding oddities without any labels}

\pw{Train a model to compress each node into a few numbers and then rebuild the
original. Because it practises on mostly normal data, it gets good at rebuilding
normal things. Anything it rebuilds \emph{badly} is unusual --- and unusual is
worth a look. Crucially this needs \textbf{no labels at all}.}

We implement two: \textbf{GAE} (rebuilds the connections) and
\textbf{DOMINANT-style} (rebuilds both connections and node attributes).

\textbf{Why does this matter?} Real deployments usually have no labels. The OCDS
Georgia dataset has none by design, and the unsupervised arm is the only thing
that works there.

\subsection{The structural floor --- deliberate simplicity}

The simplest possible scorer: count structural things (degrees, triangles,
k-core), convert to z-scores, average the positive ones. No learning at all.

\textbf{Why keep something so basic?} So we always know what ``zero cleverness''
scores. It is our reality check --- and on one important task
(\code{common\_control}) \textbf{the floor beats everything else}, which is a
result worth knowing.

\subsection{Calibration and the ensemble --- combining models fairly}

\begin{plain}
\pw{Model A says ``0.9'' and model B says ``0.4'' --- but A might be an
optimist who says 0.9 about everything. Before averaging, you must convert each
model's private scale into real probabilities. We use \emph{isotonic regression}
on validation data, which learns the correction curve without assuming any
particular shape.}
\end{plain}

\textbf{This is not a preference, it is a measured requirement.} Combining the
models \emph{without} calibrating collapses performance from
\textbf{0.4434 to 0.0511} AUC-PR. That is a catastrophic failure, and it is
seed-invariant (it happens every time). Details in Section~\ref{sec:results}.

\subsection{Leiden --- grouping into communities}

\pw{Leiden finds clusters that are densely connected inside and sparsely
connected to the outside. We run it to convert ``suspicious nodes'' into
``suspicious groups'', because investigators open cases, not rows.}

\textbf{Why Leiden and not Louvain?} Louvain --- the older, more famous
algorithm --- can produce \emph{internally disconnected} communities: it can put
two unrelated fragments in the same ``group''. For us that would mean an alert
whose members are not actually connected to each other, which would be an
incoherent case file. Leiden guarantees connected communities.

\subsection{De-duplication (NMS) and the size cap}

Go down the ranked list; if a new alert overlaps an accepted one by more than
50\% (Jaccard), suppress it. Also cap alerts at 100 members.

\pw{Without this, one big cluster would appear ten times in slightly different
forms and eat the entire review budget. And a ``community'' of 5{,}000 members
is not a lead, it is a phone book.}

\subsection{XGBoost baselines --- the honesty yardstick}

\pw{XGBoost is a classical (non-neural) method that builds many small decision
trees, each fixing the previous one's mistakes. It is fast, robust, and very
strong on table-shaped data.}

We run four baselines: \textbf{B1} hand-written rules, \textbf{B2} XGBoost on
plain features, \textbf{B3} XGBoost plus neighbour-averaged graph features (the
GADBench recipe), \textbf{B4} the classical bid-rigging screens alone.

\begin{warnbox}
\textbf{Why we deliberately built the thing most likely to beat us.} A published
benchmark (GADBench, NeurIPS 2023) found that tree models with simple
neighbourhood features often beat specialised GNNs. Any reviewer who knows that
paper will immediately ask ``did you compare against it?''. We did, and
\textbf{B3 beats our GNN on Elliptic++} --- reported openly in
Section~\ref{sec:results} rather than hidden.
\end{warnbox}

\subsection{PGExplainer --- explaining the decisions}

\pw{After the model flags a group, we ask ``which few links actually caused
this?'' The explainer learns a mask that keeps the important edges and drops the
rest, leaving a small sub-network a human can read.}

\textbf{Why PGExplainer and not GNNExplainer?} We measured both
(Section~\ref{sec:tuning}). GNNExplainer optimises separately for every alert
(slow) and produced explanations that failed a basic sanity check on
\textbf{41 of 50} alerts. PGExplainer trains once and then explains the whole
queue cheaply; sanity failures dropped to \textbf{1 of 50}, and it was
2.4$\times$ faster. Evidence-based swap.

\subsection{Why not these other methods?}

An examiner will ask this, so here are the honest answers.

\begin{longtable}{@{}p{3.9cm}p{11.6cm}@{}}
\toprule
\textbf{Not used} & \textbf{Why} \\
\midrule
\endhead
\textbf{Plain GAT} & Superseded by GATv2 at identical cost (static-attention
flaw). Kept as an ablation only. \\
\textbf{Louvain} & Can emit disconnected communities --- meaningless as a case
file. Leiden fixes exactly this. \\
\textbf{GraphSMOTE} (synthetic minority nodes) & Creating fake nodes distorts
the very topology we are trying to measure. We inject \emph{whole labelled
sub-networks} instead, which matches the real detection target. \\
\textbf{SMOTE / random oversampling} & Same objection, plus it ignores graph
structure entirely. \\
\textbf{Accuracy / ROC-AUC as headline} & Meaningless at 2\% prevalence.
ROC-AUC is optimistic under heavy imbalance; AUC-PR is the correct choice. \\
\textbf{Random train/test splits} & Would leak the future into the past.
Forbidden project-wide. \\
\textbf{TGN / TGAT} (continuous-time GNNs) & Genuinely promising, but we get
temporal behaviour through snapshots plus time encodings. Listed as future work,
not claimed. \\
\textbf{PyGOD library} & Attempted and rejected --- unusable in our setting; we
wrote native autoencoder implementations instead (recorded in the decision
log). \\
\textbf{Context-fusion gating} & \emph{Built and measured}, then \textbf{not
adopted}: it made results worse under temporal shift. A negative result we keep
and report. \\
\textbf{Line-graph view} & Same story --- built, measured ($0.4986$ vs
$0.5492$ at seed 0), not adopted on this dataset. \\
\textbf{A large language model as the detector} & LLMs do not read million-edge
graphs, cannot be calibrated to a budget, and hallucinate. The LLM is confined
to the explanation/assistant layer where every claim is checked against
retrieved evidence. \\
\bottomrule
\end{longtable}

% =====================================================================
\section{How We Measure Success}

\subsection{The alert unit --- defining a ``hit''}

Before you can measure precision on groups you must define what counts as
correct. Ours: an alert is a true positive if it contains \textbf{at least one}
confirmed illicit node / cartel firm.

\begin{warnbox}
\textbf{The obvious objection, and our answer.} A lenient rule like this
flatters models that flag enormous groups --- flag 5{,}000 nodes and you are
bound to catch someone. That is precisely why we (a) cap alerts at 100 members,
(b) de-duplicate, and (c) run a \textbf{sensitivity analysis} at stricter
thresholds (10\% and 25\% of members confirmed). The queues proved robust under
those stricter rules.
\end{warnbox}

\subsection{The metrics we report, and why}

\begin{itemize}
  \item \textbf{Precision@k} at real budgets --- the operational truth.
  \item \textbf{AUC-PR with the prevalence baseline printed beside it} --- always
        both numbers, never AUC-PR alone.
  \item \textbf{Recall@budget / FPR@budget} --- what you catch and what you waste.
  \item \textbf{Per-time-step metrics} --- so a single good period cannot hide a
        collapse elsewhere.
  \item \textbf{5 seeds, mean $\pm$ std} --- single-seed numbers are not
        publishable.
  \item \textbf{Paired bootstrap} --- to test whether a gap is real or luck.
  \item \textbf{Explanation fidelity} --- whether explanations are causally real.
\end{itemize}

% =====================================================================
\section{Tuning: What We Tried and How Results Changed}
\label{sec:tuning}

This is the iteration history, in order. It is deliberately warts-and-all,
because the corrections are the most instructive part.

\subsection{Iteration 1 --- baselines first}

Before any deep learning we built the simple methods, so we would know what
``good'' meant.

\begin{center}\small
\begin{tabular}{@{}lrr@{}}
\toprule
\textbf{Elliptic++ (prevalence 0.065)} & \textbf{AUC-PR} & \textbf{P@100} \\
\midrule
B1 rules engine & 0.0576 & 0.00 \\
B2 XGBoost & 0.8076 & 1.00 \\
B3 XGBoost + graph features & \textbf{0.8104} & 1.00 \\
\bottomrule
\end{tabular}
\end{center}

\textbf{What we learned immediately:} the hand-written rules engine scores
\emph{below} the prevalence baseline --- i.e.\ worse than random ranking. This is
the classic ``rules produce too many false positives'' complaint, measured on
real data. And the tree models are extremely strong.

\subsection{Iteration 2 --- first GNNs}

GATv2 with focal loss reached AUC-PR \textbf{0.693}, P@100 0.99. Focal loss beat
weighted cross-entropy on both GNN families. It looked like a success.

\subsection{Iteration 3 --- the audit that cut our own numbers}

We ran a deep internal audit and found \textbf{30 defects}. Two mattered enormously:

\begin{warnbox}
\textbf{Finding F3 --- accidental test-time adaptation.} The scoring pass was
re-computing normalisation statistics \emph{on the test data}. That is a subtle
form of leakage: the model was quietly absorbing part of the distribution shift
it was supposed to be measured against.

\textbf{Fix:} freeze normalisation at training time.
\textbf{Effect: GATv2 fell from 0.693 to 0.532.}

We cut our own headline by 23\% because the higher number was not honest. An
\emph{explicit} test-time-adaptation experiment is a legitimate method; an
\emph{accidental} one is a bug.
\end{warnbox}

\begin{warnbox}
\textbf{Finding --- ties in ranking.} Metrics were not tie-aware, which inflated
Precision@k for the rules and screens baselines (Mendeley B1 P@18 fell from 0.56
to 0.47 once corrected).
\end{warnbox}

Every published number was regenerated after the fixes, and 34 regression tests
were added so these specific bugs cannot return.

\subsection{Iteration 4 --- calibration before fusion}

We combined the supervised model, the autoencoders and the floor.

\begin{center}\small
\begin{tabular}{@{}lr@{}}
\toprule
\textbf{Fusion method (5 seeds)} & \textbf{AUC-PR} \\
\midrule
Naive rank fusion & 0.0511 $\pm$ 0.0019 \\
\textbf{Calibrated fusion} & \textbf{0.4434 $\pm$ 0.0501} \\
\bottomrule
\end{tabular}
\end{center}

\pw{Averaging the models raw destroys them, because the weak unsupervised
members drown out the strong supervised one. Calibrating first fixes it.}
Paired bootstrap at seed 0: $\Delta = +0.471$, 95\% CI $[0.440, 0.499]$,
$p \approx 0.001$.

\subsection{Iteration 5 --- multi-seed reality check}

Running 5 seeds instead of 1 changed the story again. The single-seed 0.5492 was
\emph{the best seed}, not a typical one.

\begin{center}\small
\begin{tabular}{@{}lrr@{}}
\toprule
\textbf{Model} & \textbf{Single seed} & \textbf{5-seed mean $\pm$ std} \\
\midrule
GATv2-focal & 0.5492 & \textbf{0.4729 $\pm$ 0.0525} \\
GATv2-weighted CE & 0.4869 & 0.4388 $\pm$ 0.0505 \\
\bottomrule
\end{tabular}
\end{center}

\textbf{Consequence:} the focal-vs-weighted-CE question, which looked settled at
one seed, is actually \textbf{second-order} --- the two overlap within their
spread. We report it as inconclusive on GATv2 rather than claiming a win.
(Note: the weighted-CE figure above is this machine's own 5-seed run,
$0.4388 \pm 0.0505$; a collaborator machine measured $0.4435 \pm 0.0615$ --- the
difference is within the recorded $\pm0.02$ machine-class variance and the
verdict is identical.)

\subsection{Iteration 6 --- explainer swap}

\begin{center}\small
\begin{tabular}{@{}lrr@{}}
\toprule
 & \textbf{GNNExplainer} & \textbf{PGExplainer} \\
\midrule
Passes PyG sanity check & 12 / 50 & \textbf{49 / 50} \\
Nonsensical fidelity & 38--41 / 50 & \textbf{1 / 50} \\
Speed (amortised) & 1$\times$ & \textbf{2.4$\times$ faster} \\
\bottomrule
\end{tabular}
\end{center}

\subsection{Iteration 7 --- component ablations}

\begin{center}\small
\begin{tabular}{@{}lrl@{}}
\toprule
\textbf{Change} & \textbf{Effect on AUC-PR} & \textbf{Verdict} \\
\midrule
Remove bidirectional edges & \textbf{$-$0.194} (seed 0) & Most important component \\
Remove unsupervised members & \textbf{$+$0.030} (5 seeds) & Removing them \emph{helps} \\
Remove focal loss & $-$0.034 (5 seeds) & Within noise --- second-order \\
Add context-fusion gating & $-$0.225 (seed 0) & Rejected \\
Add line-graph channel & $-$0.051 (seed 0) & Rejected \\
Add raw+structural channels & $-$0.171 (seed 0) & Rejected \\
Add screens to XGBoost (Mendeley) & \textbf{$+$0.063} & Adopted \\
\bottomrule
\end{tabular}
\end{center}

\begin{plain}
\textbf{Reading this table correctly.} The middle column is the effect
\emph{of making the change}. A ``remove'' row therefore measures what a
component contributes (a negative number means the component was helping). An
``add'' row measures a variant we tested and \emph{rejected}. These are
different kinds of claim and must not be compared with each other --- only
``remove'' rows support a ``most important component'' statement. Our table
builder enforces this separation automatically, and also keeps 5-seed, seed-0
and deterministic comparisons in separate blocks so their numbers can never be
subtracted from one another by accident.
\end{plain}

\textbf{The honest reading of the unsupervised arm:} taking it out
\emph{improves} the fused score by $0.030$ (from $0.4434$ to $0.4729$) --- i.e.\
including it \emph{costs} us $0.030$ AUC-PR. We keep it anyway, because it
provides \emph{coverage} where no labels exist: the entire OCDS Georgia study,
and any realistic deployment on unlabelled data, depends on it. Calibration is
what stops these weak members from destroying the fusion. We keep it for
robustness, not for the headline --- and we say so rather than quietly dropping
an arm that does not flatter us.

\subsection{Iteration 8 --- the UI, four attempts}

The dashboard was rejected by the stakeholder three times before it passed.
\textbf{V1} was rejected as monochrome and flat; \textbf{V2} added a five-hue
system, visible glass and hover states but was still called too quiet;
\textbf{V3} added real WebGL and a proper tab system and passed review \#4;
\textbf{V4} (current) answered the last complaint --- ``still blue dominated''.

We measured the V4 problem rather than guessing. The old background shader
painted \textbf{59.1\% of screen pixels a strong blue}; after the fix it is
\textbf{0\%}, with the average background colour moving from RGB(21,25,40) to
RGB(11,12,15) --- true near-black. Colour is now reserved for \emph{meaning}
(risk, category, view identity), which is how professional dashboards are built.

% =====================================================================
\section{Results}
\label{sec:results}

\subsection{Financial (Elliptic++, prevalence 0.065)}

\begin{center}\small
\begin{tabular}{@{}lrr@{}}
\toprule
\textbf{Model} & \textbf{AUC-PR} & \textbf{P@100} \\
\midrule
B1 rules & 0.0576 & 0.00 \\
B2 XGBoost & 0.8076 & 1.00 \\
\textbf{B3 XGBoost + graph} & \textbf{0.8104} & 1.00 \\
GATv2-focal (5 seeds) & 0.4729 $\pm$ 0.0525 & 0.812 $\pm$ 0.238 \\
GATv2-weighted CE (5 seeds) & 0.4388 $\pm$ 0.0505 & 0.724 $\pm$ 0.147 \\
Calibrated ensemble (5 seeds) & 0.4434 $\pm$ 0.0501 & --- \\
Rank fusion (5 seeds) & 0.0511 $\pm$ 0.0019 & --- \\
Autoencoders / floor & 0.039--0.055 & --- \\
\bottomrule
\end{tabular}
\end{center}

\begin{warnbox}
\textbf{Our headline negative result, stated plainly.} \textbf{XGBoost beats our
graph neural network on this dataset} --- 0.8104 versus 0.4729, roughly six
standard deviations apart. This replicates the GADBench finding rather than
contradicting it.

\textbf{Why?} Elliptic's 166 features \emph{already contain} 72 hand-engineered
one-hop neighbourhood aggregates. The graph information a GNN would learn has
been pre-baked into the table, so the tree model gets the structural signal for
free without the optimisation difficulty. That is also why B3 barely improves on
B2 --- adding neighbour features to features that are already neighbour features.

We report this because it is true, and because a reviewer familiar with the
literature would find it in five minutes.
\end{warnbox}

\textbf{The actor-level view} is a useful contrast: global AUC-PR is only
0.2473, yet the \emph{top} of its queue is excellent and stable (P@100 $\ge$
0.98). This is exactly why budget-first metrics matter --- a mediocre global
score can still be an excellent worklist.

\subsection{Procurement (Mendeley, prevalence 0.358, within-sample)}

\begin{center}\small
\begin{tabular}{@{}lrr@{}}
\toprule
\textbf{Model} & \textbf{AUC-PR} & \textbf{P@18} \\
\midrule
B1 rules & 0.3426 & 0.47 \\
B2 XGBoost & 0.3925 & 0.22 \\
B3 XGBoost + graph & 0.3775 & 0.56 \\
B4 screens & 0.3811 & 0.78 \\
\textbf{B2 + dataset screens} & \textbf{0.4558} & 0.61 \\
R-GCN (5 seeds) & 0.2808 $\pm$ 0.0087 & --- \\
\bottomrule
\end{tabular}
\end{center}

\begin{warnbox}
\textbf{Second honest negative.} The graph model scores \textbf{below the
prevalence baseline} (0.2808 vs 0.358) and does so \emph{consistently} across
seeds --- it is a stable negative, not noise. Causes: an era shift between the
training and test periods, and genuinely weak signal in award-only data. On this
dataset the best firm-level signal is tabular (0.4558).

We did not delete this arm. A negative result on a fair test is a finding.
\end{warnbox}

\subsection{Transfer --- does it work somewhere new?}

\begin{itemize}
  \item \textbf{García (leave-one-market-out):} positive on \emph{every} market;
        macro lift \textbf{1.57}; Italy reaches Precision@10 = 1.00.
  \item \textbf{Mendeley (leave-one-country-out):} market-dependent --- macro lift
        1.17, but \textbf{the largest market fails} (lift 0.90, i.e.\ worse than
        guessing there). An average that hides a failure is a misleading average,
        so we publish the whole matrix.
  \item \textbf{Cross-domain:} procurement$\to$financial is weakly positive
        ($0.1501$ vs $0.065$ prevalence); financial$\to$procurement is negative.
  \item \textbf{Label efficiency:} transfer \textbf{never pays} at $\le$500
        target labels in either direction. The pretrained encoder is not a
        shortcut when labels are scarce --- a clean negative answer to a question
        nobody had tested.
\end{itemize}

\subsection{Injection at scale --- the RQ2 result}

We planted 100 fake cartels (940 member firms) into the 163{,}327-node OCDS
Georgia test window and measured recovery over 5 seeds, with \emph{no labels
available} to the detectors.

\begin{center}\small
\begin{tabular}{@{}lrl@{}}
\toprule
\textbf{Planted pattern} & \textbf{Recall@2000} & \textbf{Verdict} \\
\midrule
coordinated\_cluster (clique) & \textbf{0.9225 $\pm$ 0.1733} & Found (4/5 seeds perfect) \\
common\_control & 0.4286 $\pm$ 0.0000 & Firms found deterministically \\
partition & 0.172 & Evades \\
rotation & 0.084 & Evades \\
cover\_bid & 0.010 & Evades \\
\bottomrule
\end{tabular}
\end{center}

\begin{keybox}
\textbf{The finding.} At a 1.2\% review budget, structure-only unsupervised
screening reliably catches \emph{clique-type} coordination --- firms that
physically bid together --- and is \textbf{blind} to award-pattern cartels
(rotation, market partition, cover bidding), seed-invariantly.

\pw{If the cheats leave a visible web of connections, we find them. If they
merely take turns winning without ever bidding together, structure alone cannot
see it --- you need price and sequence evidence.} This tells a deployment team
exactly which cartels this tool will and will not catch, which is far more
useful than a single average score.
\end{keybox}

\subsection{Robustness}

\begin{itemize}
  \item Queue metrics are invariant to the de-duplication threshold and robust to
        the hit rule --- the protocol is not doing the work.
  \item The OCDS ingest and the tree baselines \textbf{reproduce byte-for-byte
        across three different machines}.
  \item \textbf{Validation blindness (an important warning):} validation AUC-PR
        sits at 0.93--0.95 while test sits at 0.42--0.55. Worse, adding up to
        20\% \emph{noise} to training labels \emph{raises} test AUC-PR to
        0.5978 while validation collapses. Under temporal shift, our validation
        set is not a trustworthy guide to test performance. We state this as a
        known limitation rather than quietly using validation to pick winners.
\end{itemize}

% =====================================================================
\section{The System Architecture}
\label{sec:arch}

This section explains \code{docs/architecture.html} --- the interactive diagram
page --- in the same plain language. Open that file in any browser; it has a
light/dark toggle and needs no internet connection.

The page has four parts: the system flowchart, the algorithm cards, the AWS
cloud diagram, and the cost tables.

\subsection{The system flowchart --- ten numbered stages}

The diagram is read \textbf{top to bottom}. Everything above the trust boundary
is the \emph{offline deep-learning pipeline}; everything below is the
\emph{always-on product}.

\begin{longtable}{@{}p{0.7cm}p{4.0cm}p{10.5cm}@{}}
\toprule
\textbf{\#} & \textbf{Stage} & \textbf{What it means} \\
\midrule
\endhead
1 & Data ingestion adapters & Reads the six public datasets and checks each file
against its checksum. One adapter per dataset. \\
2 & Unified graph schema (IR) & Converts everything into the one common format,
and applies the time-safe splitting rules. After this box, nothing knows which
domain it came from. \\
3a & Supervised GNNs & GATv2, GraphSAGE, R-GCN with focal loss --- the models
that learn from labelled examples. \\
3b & Unsupervised arm & DOMINANT and GAE autoencoders --- work with no labels. \\
3c & Graph features & Degrees, motifs, k-core, bid screens, ratios --- the
hand-computed numbers. \\
3d & Baselines \& checks & The XGBoost yardstick, injection studies, and the
test suite. \\
 & & \emph{(3a--3d sit side by side because they run in parallel on the same
graph, and each produces its own score for every node.)} \\
4 & Calibrated ensemble fusion & Converts each model's scores into honest
probabilities \emph{first}, then merges them into one comparable risk score. \\
5 & Leiden $\rightarrow$ alert queue & Groups nodes into communities, ranks them,
removes duplicates, and cuts at the review budget. Produces 254 financial and
223 procurement alerts. \\
6 & Explanation layer & PGExplainer finds the minimal sub-network; the motif
matcher names the pattern; FATF/OECD citations are attached. \\
7 & Artifact store & Plain files: \code{alerts.parquet},
\code{explanations/*.json}, metrics, \code{serving.json}. Nothing else. \\
\midrule
\multicolumn{3}{@{}l@{}}{\textbf{--- TRUST BOUNDARY ---}} \\
\multicolumn{3}{@{}p{15.2cm}@{}}{\emph{Above: writes files, then exits.
Below: read-only --- cannot retrain a model or change a score.}} \\
\midrule
8 & FastAPI REST layer & A small read-only web service that serves the files:
\code{/alerts?budget=k}, \code{/subgraph/\{id\}}, \code{/explanations/\{id\}},
\code{/metrics}, \code{/rigor}. Subgraphs are trimmed server-side so the browser
never receives a million edges. Every response carries the ethics caveat. \\
-- & GenAI Investigator Copilot & Mounted alongside the API at
\code{/api/v1/copilot}. Read-only SQL and tools; every number it states must
exist in the retrieved evidence; accusatory language is auto-rewritten. \\
9 & Investigator console & The React + WebGL dashboard: ranked queue, graph
explorer, evidence dossiers, Model Lab, Copilot dock, and the financial
$\leftrightarrow$ procurement toggle. \\
10 & Investigators decide & A human makes the escalation call, at a fixed
budget. The system never decides anything. \\
\bottomrule
\end{longtable}

\begin{keybox}
\textbf{Why the trust boundary is the most important line on the diagram.}
Training is powerful and dangerous: it can change what the system believes.
Serving is what the outside world touches. By separating them completely --- the
pipeline writes files and exits; the website can only read those files --- a
compromise of the public-facing service cannot alter a single score or retrain
anything. It also means the website needs no GPU, starts instantly, and can be
demonstrated on a laptop with no internet.
\end{keybox}

\subsection{The algorithm cards}

Below the flowchart the page has explanatory cards, each with a
``\emph{What it does}'' and a ``\emph{How it works}'' line, covering: graph
neural networks, GATv2 attention, GraphSAGE and R-GCN, the autoencoders,
calibration and ensembling, Leiden grouping and de-duplication, the XGBoost
baselines, the pattern matcher with official red flags, and the statistical
rigor protocols. These are the same explanations given in
Section~\ref{sec:algos} of this report --- the page is the visual version, this
report is the reference version.

% =====================================================================
\section{The Cloud (AWS) Architecture}
\label{sec:cloud}

The second diagram on the architecture page shows how this would run as a real
hosted product. It has \textbf{two planes}, and the separation between them is
the whole point.

\subsection{Plane 1 --- the serving plane (always on, cheap, no GPU)}

\begin{longtable}{@{}p{4.3cm}p{11.2cm}@{}}
\toprule
\textbf{Component} & \textbf{What it is, in plain words} \\
\midrule
\endhead
\textbf{Analyst teams} & The users. They need only a browser --- nothing to
install. \\
\textbf{Route 53} & The phone book of the internet. Turns your domain name into
the address of the actual server. \\
\textbf{CloudFront} & A global network of caches. It keeps a copy of the
dashboard near each user so the page loads fast anywhere in the world. \\
\textbf{WAF} (Web Application Firewall) & A filter in front of everything that
blocks common web attacks before they reach us. \\
\textbf{HTTPS} & The encrypted connection, so nobody can read or tamper with the
traffic in transit. \\
\textbf{S3 --- dashboard hosting} & Cheap file storage holding the built React
website, in a private bucket that only CloudFront may read. \\
\textbf{API + Copilot compute} & The machine(s) running the read-only FastAPI
service and the Copilot. \emph{At launch:} AWS Lambda, which costs \$0 while
nobody is using it. \emph{As it grows:} Fargate containers behind a load
balancer, spread across two availability zones for resilience. \\
\textbf{S3 artifact bucket} & Versioned results storage. Every pipeline run
writes to \code{runs/<run\_id>/}. \textbf{Rolling back a bad model is just
pointing at the previous folder} --- no retraining, no panic. \\
\textbf{ECR} & The registry that stores our container images (keeps the last 3). \\
\textbf{Secrets Manager} & Where the LLM API key lives, so no key is ever in the
code or in a file. \\
\textbf{IAM (read-only)} & Permissions. The serving side is granted
\emph{read-only} rights to the artifact bucket --- the trust boundary enforced by
the cloud itself, not merely by our code. \\
\textbf{CloudWatch} & Logs, metrics and alarms --- how you find out something
broke without a user telling you. \\
\textbf{AWS Budgets} & A spending alarm, switched on from day one so a runaway
job cannot quietly cost a fortune. \\
\bottomrule
\end{longtable}

\subsection{Plane 2 --- the scheduled deep-learning plane (exists only while it runs)}

\begin{longtable}{@{}p{4.3cm}p{11.2cm}@{}}
\toprule
\textbf{Component} & \textbf{What it is, in plain words} \\
\midrule
\endhead
\textbf{EventBridge schedule} & An alarm clock. ``Every night at 2am'' or
``every Sunday'' it starts the pipeline. If the job fails, it raises an SNS
alert. \\
\textbf{Spot GPU batch job} & A rented graphics machine, at roughly 70\% off
because you accept it can be reclaimed. It runs
ingest $\rightarrow$ train $\rightarrow$ score $\rightarrow$ group
$\rightarrow$ explain, writes its results to the artifact bucket,
\textbf{and then shuts itself down}. \\
\textbf{The upward arrow} & The only connection between the planes: the batch job
\emph{writes} new results into the versioned bucket that the serving plane
\emph{reads}. There is no path in the other direction. \\
\bottomrule
\end{longtable}

\begin{keybox}
\textbf{Why this shape saves most of the money.} GPUs are the expensive part.
Naively you would keep a GPU server running all the time. But our system only
needs a GPU while \emph{training}, which is a few hours a week --- scoring and
serving need none. So the GPU exists for two hours and disappears, while the
always-on part is a tiny CPU service that costs almost nothing. This is why the
demo tier is about \$10/month and effectively \$0 on standard academic credits.
\end{keybox}

\subsection{What it would cost}

The page carries three costed tiers (list prices; the demo tier is effectively
free under the \$100--200 new-account credits).

\begin{center}\small
\begin{tabular}{@{}p{4.4cm}p{7.4cm}r@{}}
\toprule
\textbf{Tier} & \textbf{What you get} & \textbf{Approx.\ / month} \\
\midrule
1 --- University demo & S3 + CloudFront (free tier), one small EC2 t3.micro for
API + Copilot, Route 53; training on your own PC or free Colab &
$\approx$\,\$10 \\
2 --- Product launch & S3 + CloudFront, Lambda + API Gateway for up to 1M
requests, S3 + ECR, a rented GPU $\approx$4\,h/week, pay-per-use LLM,
monitoring and secrets & $\approx$\,\$14--55 \\
3 --- Growing product & Same plus ECS Fargate across two zones, load balancer,
WAF, daily GPU refresh $\approx$2\,h/day, team-scale LLM usage &
$\approx$\,\$90--183 \\
\bottomrule
\end{tabular}
\end{center}

% =====================================================================
\section{The Dashboard --- Every Screen Explained}
\label{sec:dashboard}

\subsection{Starting it}

\begin{plain}
\textbf{With real results:}\\
\code{uv run collusiongraph serve --port 8001}\\
then, in \code{frontend/}: \code{VITE\_API\_TARGET=http://127.0.0.1:8001 npm run dev}\\
open \code{http://localhost:5173}

\vspace{4pt}
\textbf{Without any trained results} (any laptop, no API keys):\\
\code{uv run python scripts/build\_dev\_store.py} then \code{poe serve}

\vspace{4pt}
\textbf{Everything at once:} \code{docker compose up} $\rightarrow$ port 8080
\end{plain}

\subsection{Things that appear on every screen}

\begin{itemize}
  \item \textbf{Domain toggle} --- switches financial $\leftrightarrow$
        procurement. The accent colour changes with it.
  \item \textbf{Dataset selector} --- two datasets per domain
        (\code{elliptic\_pp} / \code{elliptic\_pp\_actor}; \code{mendeley\_eu} /
        \code{garcia\_rodriguez}).
  \item \textbf{Six tabs}, each with its own fixed colour and icon.
  \item \textbf{The ethics caveat}, in the footer of every screen, always.
  \item \textbf{Copilot dock} --- the assistant, openable from anywhere.
\end{itemize}

\subsection{Screen 1 --- Overview (command deck)}

\begin{longtable}{@{}p{4.3cm}p{11.2cm}@{}}
\toprule
\textbf{Element} & \textbf{What it means} \\
\midrule
\endhead
\textbf{Alerts in queue} & Total suspicious groups found in this dataset. \\
\textbf{High-band alerts} & How many are in the top risk band. \\
\textbf{Budget k slider} & \emph{The most important control on the page.} How
many cases you can review. Everything re-ranks live as you drag it. \\
\textbf{Precision@k readout} & The measured share of the top $k$ that are real
--- it updates as you move the slider, which makes the central metric of the
whole project tangible. \\
\textbf{Flagged-community constellation} & The hero visual: each dot is an alert,
size = number of members, colour = risk band (high / medium / low). Layout is
schematic but \emph{ranks are real}. \\
\textbf{Top flagged communities} & The leaderboard: rank, risk score, member
count, jump arrow. \\
\bottomrule
\end{longtable}

\subsection{Screen 2 --- Alert Queue (the working surface)}

The ranked table an investigator actually works through. Each row shows:

\begin{itemize}
  \item \textbf{Rank} and \textbf{calibrated risk score} (a real probability,
        because of the calibration step --- not an arbitrary 0--1 number).
  \item \textbf{Motif chip} --- the named pattern, with its own glyph and colour.
  \item \textbf{Community size} --- how many entities are involved.
  \item \textbf{Red-flag count} --- how many official indicators matched.
  \item \textbf{Activity sparkline} --- a tiny chart of activity over time. A flat
        line means all the activity sits in one time window (common in the
        Elliptic data); a varied line means genuine spread.
  \item \textbf{Row actions} --- open the graph, open the dossier, ask the Copilot.
\end{itemize}

\subsection{Screen 3 --- Graph Explorer}

A WebGL canvas showing the alert's network.

\begin{longtable}{@{}p{4.3cm}p{11.2cm}@{}}
\toprule
\textbf{Element} & \textbf{What it means} \\
\midrule
\endhead
\textbf{Coral nodes} & Flagged members of this alert. Coral is used for nothing
else anywhere in the product. \\
\textbf{Grey nodes} & Context --- neighbours shown so the group is readable. \\
\textbf{Lit edges} & Connections currently revealed by the replay. \\
\textbf{Replay control} & Animates the network being built up. Two honest modes:
\textbf{``replay flow''} when the edges carry different timestamps (real time
order), and \textbf{``replay order''} when they do not (sequence order). The
readout tells you which one you are watching --- e.g.\ ``all edges $\cdot$ window
38'' versus ``full window 2014--2018''. \\
\textbf{Click a node} & Expands its neighbours. \\
\bottomrule
\end{longtable}

\begin{plain}
\textbf{Why two replay modes exist} --- a real bug we found and fixed. Elliptic
transaction edges within one alert often share a \emph{single} time step (we
measured: the top alert's 101 edges all sit at window 38). The original replay
animated by timestamp, so with all timestamps equal nothing appeared to move and
the control looked broken. Rather than fake a time spread, we detect the
situation and switch to sequence order, and we \emph{say so on screen}.
\end{plain}

\subsection{Screen 4 --- Case Detail (the evidence dossier)}

Where an investigator decides. Contents:

\begin{itemize}
  \item \textbf{Motif schematic} --- an animated drawing of the detected shape.
  \item \textbf{Minimal subgraph} --- the few links the explainer says actually
        caused the score.
  \item \textbf{Evidence fields} --- amounts and fees where the dataset has them,
        structural and temporal evidence where it does not. \emph{We never invent
        a field the data cannot support.}
  \item \textbf{Red-flag cards} --- e.g.\ \code{FATF-STRUCT-01} ``sub-threshold
        deposits converging on one account'', or \code{OECD-ROT-01} ``the same
        supplier group wins in turn'', each with \emph{why it matched}.
  \item \textbf{Evidence-source labels} --- each item marked as coming from the
        learned model, the structural rules, or the statistical screens. This
        prevents the false impression that one method explains everything.
  \item \textbf{Fidelity scores} --- how much the score depends on the highlighted
        evidence.
  \item \textbf{Export} --- download the alert as JSON.
  \item \textbf{The locked caveat} --- screening signal only. A code-level
        validator refuses to construct a bundle with altered wording.
\end{itemize}

\subsection{Screen 5 --- Model Lab}

The figure factory: PR curves, precision-versus-budget curves with markers, the
transfer matrix heat map (train market $\times$ test market), ablation
comparisons, and the rigor panel (multi-seed spreads, confidence intervals,
label-noise curve). Everything exports to SVG/PNG for papers and slides.

\subsection{Screen 6 --- About}

A scroll-driven explanation of the ``two ledgers, one structure'' idea, with the
motif table animated. Doubles as the opening narrative for a demo.

\subsection{The Copilot dock}

Ask questions in ordinary English (``why was alert 3 flagged?''). What makes it
safe:

\begin{itemize}
  \item \textbf{Read-only.} It can query and read; it cannot change anything.
  \item \textbf{Grounding gate.} Every number it states must exist in the
        evidence it retrieved. Ungrounded claims are blocked, not printed.
  \item \textbf{Guilt-language guard.} Accusatory phrasing is rewritten or
        blocked before release. Measured: \textbf{zero released violations}
        across the golden-question gate.
  \item \textbf{24 golden questions} run as a gate; all 24 pass.
  \item \textbf{Every answer is labelled AI-generated} and carries the caveat.
\end{itemize}

\textbf{Two measured weaknesses} (we tested it adversarially and record what
failed): (a) correct arithmetic \emph{derived} from grounded numbers is
sometimes flagged as ungrounded, lowering confidence unnecessarily; (b) very
open-ended list questions can exhaust the tool-call budget and return ``please
narrow the question''. Both are recorded as open work items.

% =====================================================================
\section{Who Would Actually Use This}

\subsection{A realistic day in the life}

\begin{plain}
A procurement-integrity analyst at a national audit office can review about
20 cases a week. There were 40{,}000 tenders last quarter.

\textbf{Without this system:} she filters by simple rules --- single-bid tenders,
unusual values. The list is thousands long, mostly boring, and rule-based
filters are provably poor at ranking (our B1 baseline scored \emph{below} random
on Elliptic).

\textbf{With this system:} she sets the budget slider to 20. She receives 20
ranked \emph{groups}, each named with a pattern, each carrying an OECD citation
and the specific evidence. She opens \#1: five firms, a rotation pattern, they
have never bid against each other, two share a registered address. She reads the
dossier, asks the Copilot to summarise the timeline, exports the JSON, and
escalates to a human investigator.

\textbf{She still makes every decision.} The system changed \emph{what she looks
at first}, not \emph{what is true}.
\end{plain}

\subsection{Who can use it}

\begin{longtable}{@{}p{4.3cm}p{11.2cm}@{}}
\toprule
\textbf{User} & \textbf{Value} \\
\midrule
\endhead
Public procurement / audit offices & Rank tenders and suppliers for review under
a fixed staffing budget; citations map to OECD guidance they already use. \\
Financial intelligence units, bank AML teams & Prioritise networks rather than
single transactions; FATF indicator citations map to existing reporting
vocabulary. \\
Competition authorities & Screen for bid-rigging structures across markets;
cross-market transfer results say honestly where it will and will not work. \\
Development banks / donors & Screen contracts in countries with no labelled
cartel history --- exactly the LOCO/LOMO scenario we measured. \\
Auditors and investigative journalists & Open OCDS data plus this pipeline finds
structures worth asking about. \\
Researchers & A reproducible, leakage-safe benchmark harness with honest
baselines; every number regenerates from a config. \\
\bottomrule
\end{longtable}

\subsection{Who should \emph{not} use it, and for what}

\begin{warnbox}
\textbf{Not for:} deciding guilt; automated adverse action of any kind
(blacklisting, freezing, debarment); evidence in court; producing any output
about a named private individual; real-time transaction blocking.

This is a \textbf{prioritisation aid} at TRL 3--4. It has no authentication, no
multi-user hardening, no case-management workflow, and no real-time scoring.
Fraud screening is an EU AI Act \emph{high-risk} category with obligations
enforceable from 2 August 2026; our human-in-the-loop framing, explanation
bundles and logged caveats are our documented alignment posture --- not a
compliance certification.
\end{warnbox}

% =====================================================================
\section{Limitations --- Stated Honestly}
\label{sec:limits}

\begin{enumerate}
  \item \textbf{XGBoost beats our GNN on Elliptic++} (0.8104 vs 0.4729). The
        deep model is not the best detector on our main financial dataset. Its
        value is that it produces communities and explanations; but the honest
        headline is that a tree wins on ranking.
  \item \textbf{The R-GCN scores below prevalence on Mendeley} (0.2808 vs
        0.358) --- a stable negative.
  \item \textbf{Validation blindness under temporal shift.} Validation says
        0.93--0.95, test says 0.42--0.55; label noise \emph{improves} test while
        validation collapses. We cannot currently trust validation for model
        selection under shift.
  \item \textbf{Seed variance ($\pm$0.05) is large} --- larger than most effects
        we would like to claim. It dominates machine-class variance ($\pm$0.02).
  \item \textbf{Mendeley prevalences are case-control artifacts} --- only lift
        comparisons are meaningful.
  \item \textbf{Elliptic features are anonymised}, so financial explanations lean
        structural/temporal rather than semantic.
  \item \textbf{No learned explanation for the R-GCN} --- mask-based explainers
        work only on full-edge-set convolutions; heterogeneous explanation is
        outstanding.
  \item \textbf{Human validation is pending.} The practitioner study is built and
        tested but the human phase has not been run, so ``are these explanations
        useful to real investigators?'' is not yet answered empirically.
  \item \textbf{Structure-only screening is blind to award-pattern cartels}
        (rotation, partition, cover bidding) --- measured, seed-invariant.
  \item \textbf{AMLworld is not fully exercised}, and the largest-scale GPU
        experiments (PNA, GIN+EU) remain gated on hardware we do not have.
  \item \textbf{Not a production system:} no auth, no multi-tenancy, no
        real-time path.
\end{enumerate}

% =====================================================================
\section{Future Work}
\label{sec:future}

\subsection{Directly addressing the limitations}

\begin{enumerate}
  \item \textbf{Shift-aware model selection.} The validation-blindness result is
        the single most important open problem: build selection criteria that
        survive temporal shift (e.g.\ time-sliced validation, or selecting on a
        held-out \emph{late} slice).
  \item \textbf{Close the tree-vs-GNN gap honestly.} Either find an architecture
        that genuinely beats B3 on Elliptic++, or characterise precisely when
        graph learning pays --- our own evidence suggests it pays when
        neighbourhood aggregates are \emph{not} already baked into the features.
  \item \textbf{Heterogeneous explanations} for R-GCN, closing the procurement
        explanation gap.
  \item \textbf{Run the practitioner study} with $\ge$5 raters, and report
        inter-rater agreement.
  \item \textbf{Price- and sequence-aware detectors} to attack the cartels that
        structure alone provably misses --- the clearest scientific gap we found.
  \item \textbf{GPU-scale experiments:} PNA and GIN+EU on AMLworld for external
        comparability; neighbour-sampling minibatches for graphs beyond memory.
  \item \textbf{Temporal GNNs (TGN/TGAT)} to test whether continuous-time message
        passing beats snapshots plus time encodings.
  \item \textbf{Firm--firm co-bid projection before grouping} in García, which
        would yield a richer procurement queue than the current honest-but-thin
        3 alerts.
  \item \textbf{Production hardening:} authentication, multi-tenancy, incremental
        scoring, case-management integration.
\end{enumerate}

\subsection{Should we add Web3 / blockchain? An honest assessment}
\label{sec:web3}

The question deserves a real answer rather than an enthusiastic one. We
evaluated it against a single test: \emph{does it solve a problem we actually
have?}

\begin{keybox}
\textbf{First, a fact that is easy to miss:} this project is \textbf{already}
on-chain analytics. Elliptic/Elliptic++ \emph{is} Bitcoin transaction data. We
are already doing graph forensics on a blockchain ledger. So the question is not
``should we touch Web3'', it is ``which \emph{additional} Web3 capabilities earn
their place''.
\end{keybox}

\subsubsection*{Where it genuinely adds value}

\begin{enumerate}
  \item \textbf{Live on-chain ingestion (strongest case).} Today the financial
        arm is retrospective: a frozen 2016-era Bitcoin snapshot. A direct node
        or indexer feed would make it \emph{operational} --- scoring wallets as
        activity happens. This is a genuine capability gain, and the architecture
        already supports it: it is a new adapter (stage 1) feeding the same IR.
        The honest catch is that our current pipeline is batch by design, so this
        also requires incremental scoring.
  \item \textbf{Cross-chain and bridge analysis.} Moving assets across chains via
        bridges is now a primary laundering route, and it is \emph{exactly} a
        graph problem: bridges are edges between sub-graphs. Our IR already
        supports multiple relation types and R-GCN already handles typed edges,
        so a \code{bridges\_to} edge type is a natural extension rather than a
        bolt-on. Genuinely valuable.
  \item \textbf{DeFi-native motif families.} Mixer usage (e.g.\ Tornado-style
        pools), wash trading in NFT markets, and circular flash-loan structures
        are new \emph{shapes} --- and shapes are what this system detects. They
        would slot into the existing motif table, injector and matcher without
        architectural change. This is the most natural fit of all: it extends the
        thing the project is actually good at.
  \item \textbf{Smart-contract-call graphs.} Contract-to-contract call networks
        are a rich, under-explored graph domain, and would be a genuine third
        domain testing the ``one stack, many crimes'' thesis.
\end{enumerate}

\subsubsection*{Where it would be decoration --- and one place it would be actively harmful}

\begin{enumerate}
  \item \textbf{Putting alerts or scores on a blockchain: no, and we should say
        so firmly.} It conflicts directly with the project's central ethical
        rule. Our alerts are \emph{screening signals, not findings of guilt}, and
        they can be wrong --- our own measurements say so. A blockchain is
        immutable and public by design. Publishing ``firm X is flagged'' to an
        immutable public ledger would create a permanent, unretractable
        accusation against an entity that has been convicted of nothing. That is
        the precise opposite of what \code{docs/ethics\_and\_scope.md} requires,
        and no efficiency argument outweighs it.
  \item \textbf{Tokenised incentives for reporting suspicious activity: no.} Pay
        people per flag and you manufacture false positives. Our whole design
        fights false positives under a fixed budget; a token reward would fight
        the design.
  \item \textbf{Decentralised / federated on-chain training: not justified
        here.} Federated learning solves \emph{data cannot leave the
        institution}. Every dataset we use is public. We would pay a large
        complexity and coordination cost to solve a problem we do not have.
        (It \emph{would} become relevant if real banks ever pooled private data
        --- but then the technology needed is federated learning with secure
        aggregation, which does not require a blockchain.)
  \item \textbf{NFTs, DAOs, or a project token: no.} No mechanism here needs
        transferable ownership or on-chain governance.
  \item \textbf{Blockchain as a tamper-evident audit log: technically valid, but
        not the right tool.} Auditability of ``which model version produced this
        alert'' is a real requirement. But we already meet it more simply:
        versioned artifact folders, pinned configs, checksummed inputs, and
        byte-reproducibility verified across three machines. If stronger
        guarantees were needed, a signed append-only transparency log gives the
        same tamper-evidence without a chain's cost, latency and governance
        overhead.
\end{enumerate}

\begin{keybox}
\textbf{Recommendation.} Pursue Web3 \emph{as a data source}, not as
infrastructure. Items 1--4 above (live ingestion, bridges, DeFi motifs, contract
call graphs) are worth real effort because they feed the graph engine we
already have. Everything in the second list should be declined explicitly --- and
the alerts-on-chain idea should be declined on \emph{ethical} grounds, which is
a stronger and more defensible position than declining it on technical ones.
\end{keybox}

% =====================================================================
\section{Logs, Testing and Reproducibility}

\subsection{The test suite}

\begin{center}\small
\begin{tabular}{@{}lr@{}}
\toprule
\textbf{Suite} & \textbf{Status} \\
\midrule
Backend (pytest) & \textbf{376 passing} \\
Frontend (vitest) & \textbf{35 passing} \\
Leakage suite (subset, runs standalone in CI) & 25 tests, green \\
Copilot golden questions & 24 / 24 \\
Linting (ruff), formatting (black), typing (mypy) & clean \\
Continuous integration & green on every push \\
\bottomrule
\end{tabular}
\end{center}

Notable test categories: leakage negative controls; the flagship test that the
motif matcher recovers \textbf{all ten injected families at 100\% recall} on
fixtures (matcher and injector are independent implementations, so they
cross-validate each other); calibration monotonicity; overfit-single-batch
sanity; the ethics-caveat lock; and a bidirectional check that every experiment
config appears in the reproducibility map and vice versa.

\subsection{What the pipeline writes}

Every run writes a machine-readable record --- these \emph{are} the logs:

\begin{itemize}
  \item \code{run.json} --- full configuration, seed, epochs run, best epoch, best
        validation score, training seconds, and all resulting metrics.
  \item \code{metrics.json} --- evaluation results including per-time-step
        breakdowns.
  \item \code{scores\_test.parquet} / \code{scores\_val.parquet} --- raw per-node
        scores.
  \item \code{model.pt}, \code{feature\_stats.json} --- weights and the
        \emph{frozen} normalisation constants.
  \item \code{multiseed.json} --- per-seed values plus mean/std.
  \item \code{alerts.parquet}, \code{explanations/*.json} --- the queue and the
        bundles.
  \item \code{BUILD\_REPORT.json} --- which paper tables were built, and for any
        that were skipped, the exact missing file path.
\end{itemize}

\begin{plain}
\textbf{Example of a real log record} (Elliptic++ GATv2-focal, seed 0):
model \code{gatv2}, features \code{raw}, loss \code{focal} ($\gamma=2.0$),
hidden 64, 2 layers, 4 heads, dropout 0.3, lr 0.005 --- \textbf{142 epochs run,
best at epoch 121}, best validation AUC-PR 0.9497, training time 2194.6 s,
final test AUC-PR 0.5492 at prevalence 0.0650.

\textbf{Notice the gap} between validation (0.9497) and test (0.5492) --- that
single log line \emph{is} the validation-blindness limitation, visible in the
raw record.
\end{plain}

\subsection{Reproducibility}

\begin{itemize}
  \item \code{docs/reproducibility.md} maps \textbf{all 40 experiment configs} to
        the exact headline numbers they produce. A test enforces this
        bidirectionally --- a config missing from the map, or a phantom config in
        the map, fails CI.
  \item Environment is frozen (\code{uv.lock}, pinned Python).
  \item \code{uv run poe paper-tables} regenerates every paper table from stored
        artifacts. A table whose source artifact is missing is \textbf{skipped
        with the missing path recorded} --- never rendered from partial data.
        \textbf{Current status on this machine: all 10 tables build, zero skips.}
  \item \textbf{Verified cross-machine:} the tree baselines, the LOCO fold
        (0.8025340470101002 --- to the last digit), the García matrix and the OCDS
        ingest (451{,}346 $\rightarrow$ 488{,}300 / 1{,}449{,}077) reproduce
        \emph{byte-identically} on three different machines. Torch-based results
        reproduce same-machine and shift within $\pm$0.02 across machine classes,
        which is smaller than the $\pm$0.05 seed spread --- so multi-seed means
        are the only numbers we publish.
\end{itemize}

\subsection{Adversarial self-review}

We ran a formal internal red-team pass against our own work
(\code{docs/red\_team\_review.md}), checking leakage, imbalance reporting,
baseline fairness, transfer honesty, number consistency, ablation completeness,
explanation validity and ethics propagation.

\textbf{Result: no finding invalidated any published number or protocol
guarantee.} Two presentation defects were found and fixed --- most notably
\textbf{RT-1}: single-seed bootstrap differences were printed next to multi-seed
averages, so a reader recomputing the difference from the averages would land
outside the stated confidence interval. Both are now labelled explicitly, and
our table builder enforces the separation structurally.

% =====================================================================
\section{Questions an Examiner Will Ask}

\begin{longtable}{@{}p{5.0cm}p{10.5cm}@{}}
\toprule
\textbf{Question} & \textbf{Answer} \\
\midrule
\endhead
\emph{Your GNN loses to XGBoost. Why is this a deep-learning project?} &
Because we report it. The GNN produces the community structure, the minimal
sub-network and the attention evidence that a tree cannot; and on Elliptic++
specifically the trees win because 72 neighbourhood aggregates are already baked
into the features. We state the condition under which graph learning pays rather
than claiming a win we did not get. \\
\addlinespace
\emph{How do I know you have no leakage?} & Four independent mechanisms:
strict-inductive splitting (test-period \emph{edges} withheld --- verified 77{,}512
on Elliptic++), as-of features (verified exactly equal to the truncated graph),
a 25-test leakage suite with negative controls in CI, and an audit that
\emph{found} a leak (F3) and cut our headline by 23\% to fix it. \\
\addlinespace
\emph{Is your alert-level hit rule too lenient?} & Possibly on its own, which is
why it is paired with a 100-member cap, Jaccard de-duplication and a sensitivity
analysis at 10\% and 25\% thresholds. Queue metrics were robust. \\
\addlinespace
\emph{Why AUC-PR and not ROC-AUC?} & At 2\% prevalence ROC-AUC is misleadingly
optimistic. AUC-PR's baseline is the prevalence, which we print beside every
value. \\
\addlinespace
\emph{Your baselines got the same tuning budget?} & Yes --- and precisely: no
hyperparameter search was performed on \emph{either} side. That is conservative
in our disfavour, since the trees win every comparison they enter. Stated
explicitly. \\
\addlinespace
\emph{Are the differences statistically significant?} & Paired stratified
bootstrap, 2{,}000 resamples, over identical test nodes. Reported with CIs --- and
labelled as seed-0 comparisons wherever they sit beside multi-seed means. \\
\addlinespace
\emph{One seed?} & Never for a headline. Five seeds, mean $\pm$ std. The
single-seed 0.5492 was the best seed; the mean is 0.4729. \\
\addlinespace
\emph{Are the explanations real or decorative?} & Measured with fidelity, which
is why we \emph{replaced} GNNExplainer: it failed sanity on 41/50 alerts.
PGExplainer fails 1/50. Human validation is still pending and we say so. \\
\addlinespace
\emph{Does it transfer?} & Partly, and we publish the failures: García positive
on every market (lift 1.57); Mendeley market-dependent with its largest market
\emph{failing} (0.90); cross-domain asymmetric and weak; and transfer never pays
below 500 target labels in either direction. \\
\addlinespace
\emph{What can it not catch?} & Rotation, market partition and cover bidding, at
scale, seed-invariantly --- because structure alone cannot see them. Measured on
100 planted cartels in a 163k-node network. \\
\addlinespace
\emph{Could this harm someone?} & The rule is enforced in code, not just
asserted: an immutable caveat field a validator refuses to alter, the caveat on
every API response and screen, and a Copilot guard that rewrites accusatory
language (zero released violations). No per-person outputs exist anywhere. \\
\bottomrule
\end{longtable}

% =====================================================================
\section{Summary}

\begin{keybox}
CollusionGraph turns two different crimes into the same question --- \emph{which
connected group is behaving like a coordinated ring?} --- and answers it with a
ranked, explained, budget-aware queue that a human reviews.

\textbf{What works:} clique-type coordination is found reliably at scale with no
labels (0.9225 recall at a 1.2\% budget); calibrated fusion is a measured
requirement, not a preference; transfer to unseen markets works on one dataset
family and is honestly market-dependent on another; explanations are validated
by fidelity; the whole thing reproduces byte-for-byte across machines.

\textbf{What does not:} a tree model beats our GNN on the main financial
dataset; the procurement graph model scores below its prevalence baseline;
validation cannot be trusted under temporal shift; and structure alone is blind
to award-pattern cartels.

Reporting both lists is the point. A screening tool whose limits are unknown is
not usable by anyone who is accountable for the decision.
\end{keybox}

\vspace{1em}
\begin{plain}
\textbf{Where to look next}\\
\code{PROGRESS.md} --- the full project ledger, every result with its date and commit\\
\code{docs/model\_card.md} --- the formal model card\\
\code{docs/reproducibility.md} --- all 40 configs mapped to their numbers\\
\code{docs/red\_team\_review.md} --- our adversarial self-review\\
\code{docs/architecture.html} --- the interactive architecture diagrams (Sections~\ref{sec:arch}, \ref{sec:cloud})\\
\code{docs/ethics\_and\_scope.md} --- the binding ethical constraints\\
\code{docs/DATASETS.md} and \code{docs/datasheets/} --- per-dataset detail
\end{plain}

\end{document}
